%%%%%%%%%%%%%%%%%%%%%%%%%%%%%%%%%%%%%%%%%
% fphw Assignment
% LaTeX Template
% Version 1.0 (27/04/2019)
%
% This template originates from:
% https://www.LaTeXTemplates.com
%
% Authors:
% Class by Felipe Portales-Oliva (f.portales.oliva@gmail.com) with template 
% content and modifications by Vel (vel@LaTeXTemplates.com)
%
% Template (this file) License:
% CC BY-NC-SA 3.0 (http://creativecommons.org/licenses/by-nc-sa/3.0/)
%
%%%%%%%%%%%%%%%%%%%%%%%%%%%%%%%%%%%%%%%%%

%----------------------------------------------------------------------------------------
%	PACKAGES AND OTHER DOCUMENT CONFIGURATIONS
%----------------------------------------------------------------------------------------

\documentclass[
  12pt,
  %fleqn, % Default font size, values between 10pt-12pt are allowed
  %letterpaper, % Uncomment for US letter paper size
  %spanish, % Uncomment for Spanish
]{fphw}

% Template-specific packages
\usepackage[utf8]{inputenc} % Required for inputting international characters
\usepackage[T1]{fontenc} % Output font encoding for international characters
\usepackage{mathpazo} % Use the Palatino font

\usepackage{graphicx} % Required for including images

\usepackage{booktabs} % Required for better horizontal rules in tables

\usepackage{listings} % Required for insertion of code

\usepackage{enumerate} % To modify the enumerate environment

\usepackage{enumitem}

\usepackage{amsmath}

\usepackage{mathtools,amsthm}

\usepackage{float}

\usepackage{listings}

\usepackage{hyperref}

\lstdefinelanguage{R}{
  keywords={if, else, repeat, while, function, for, in, next, break},
  otherkeywords={TRUE, FALSE, NULL, NA, Inf, NaN},
  sensitive=true,
  morecomment=[l]\#,
  morestring=[b]",
}

\newcommand{\E}{\mathbb{E}}
\newcommand{\Var}{\operatorname{Var}}
\newcommand{\R}{\mathbb{R}}
\newcommand{\norm}[1]{\left\lVert #1 \right\rVert}


%----------------------------------------------------------------------------------------
%	ASSIGNMENT INFORMATION
%----------------------------------------------------------------------------------------

\title{Homework \#6} % Assignment title

\author{Shizhe Zhang} % Student name

\date{\today} % Due date

\institute{University of California, Berkeley \\ Department of Statistics} % Institute or school name

\class{EECS 227AT} % Course or class name

\professor{Gireeja Ranade} % Professor or teacher in charge of the assignment

%----------------------------------------------------------------------------------------

\begin{document}

\maketitle % Output the assignment title, created automatically using the information in the custom commands above

%----------------------------------------------------------------------------------------
%	ASSIGNMENT CONTENT
%----------------------------------------------------------------------------------------

\section*{Problem 1. Best Approximation}
We start by recalling the Eckart-Young Theorem: Consider any square matrix $C \in \mathbb{R}^{n\times n}$ and assume that we
can write its full singular value decomposition as $C = U\Sigma V^\top$ where $\Sigma$ is the $n\times n$ diagonal matrix with distinct
diagonal entries $\sigma_1 > \sigma_2 > \cdots > \sigma_n > 0$. Then, for $0 \le k \le n$, the Eckart-Young Theorem for Frobenius
norm states that
\[
C_k = U_k \Sigma_k V_k^\top
= \underset{B \in \mathbb{R}^{n\times n}, \ \operatorname{rank}(B)\le k}{\operatorname{argmin}} \|C - B\|_F
\]
where $U_k \in \mathbb{R}^{n\times k}$ is the matrix consisting of the first $k$ columns of $U$, 
$V_k \in \mathbb{R}^{n\times k}$ is the matrix consisting of the first $k$ columns of $V$, 
and $\Sigma_k$ is the $k\times k$ diagonal matrix with the top-$k$ singular values 
$\sigma_1 > \cdots > \sigma_k$ as its diagonal entries.

(a) Now, consider any square matrix $C \in \mathbb{R}^{n\times n}$ and let $C_k \in \mathbb{R}^{n\times n}$ denote its best rank-$k$ approximation in
the Frobenius norm. Then, for any orthonormal matrix $W \in \mathbb{R}^{n\times n}$, show that
\[
W C_k W^\top
=
\arg\min_{B \in \mathbb{R}^{n\times n}, \ \operatorname{rank}(B)\le k}
\|W C W^\top - B\|_F .
\]

For the remainder of this problem, we consider a square matrix $A \in \mathbb{R}^{n\times n}$ and assume $A$ has full rank. Using
Gram-Schmidt Orthonormalization (GSO), we can write the matrix $A$ as
\[
A = QR \tag{1}
\]
where $Q \in \mathbb{R}^{n\times n}$ is an orthonormal matrix and $R \in \mathbb{R}^{n\times n}$ is an upper triangular matrix.

(b) Find the best rank-$n$ approximation to $AA^\top$ in the Frobenius norm. Justify your answer.

(c) Recall that $A \in \mathbb{R}^{n\times n}$ is a square matrix with full rank, and we can write the matrix $A$ as
\[
A = QR \tag{2}
\]
where $Q \in \mathbb{R}^{n\times n}$ is an orthonormal matrix and $R \in \mathbb{R}^{n\times n}$ is an upper triangular matrix.

Assume that $R = \operatorname{diag}(r_1, r_2, \dots, r_n) \in \mathbb{R}^{n\times n}$ is a diagonal matrix with
\[
|r_1| > |r_2| > \cdots > |r_n|,
\]
and all $r_1, r_2, \dots, r_n$ are real numbers. Let $k < n$. Then, show that the best rank-$k$ approximation to $AA^\top$
in the Frobenius norm is $Q S Q^\top$, where $S$ is a diagonal matrix defined as
\[
S = \operatorname{diag}(r_1^2, \dots, r_k^2, 0, \dots, 0) \in \mathbb{R}^{n\times n}.
\]

(d) Recall that $A \in \mathbb{R}^{n\times n}$ is a square matrix with full rank, and we can write the matrix $A$ as
\[
A = QR \tag{3}
\]
where $Q \in \mathbb{R}^{n\times n}$ is an orthonormal matrix and $R \in \mathbb{R}^{n\times n}$ is an upper triangular matrix.

Now, we no longer assume that $R$ is diagonal. Let $k < n$ and assume that the best rank-$k$ approximation to
$RR^\top \in \mathbb{R}^{n\times n}$ in the Frobenius norm is given by $G \in \mathbb{R}^{n\times n}$. Then, using the result of part (a), show that
the best rank-$k$ approximation to $AA^\top$ in the Frobenius norm is given by $Q G Q^\top$.


\subsection*{Answer}

\begin{enumerate}
\item[(a)] Let $W\in\mathbb{R}^{n\times n}$ be orthonormal, i.e. $W^\top W = WW^\top = I$.
Consider the optimization problem
\[
\min_{B:\ \operatorname{rank}(B)\le k}\ \|WCW^\top - B\|_F.
\]
Define the change of variables
\[
\widetilde{B} \doteq W^\top B W.
\]
Then $\operatorname{rank}(\widetilde{B})=\operatorname{rank}(B)$ because multiplication by an invertible matrix
(on the left or right) does not change rank. Hence $\operatorname{rank}(B)\le k$ iff $\operatorname{rank}(\widetilde{B})\le k$.

Moreover, using unitary (orthonormal) invariance of the Frobenius norm,
\[
\|WCW^\top - B\|_F
=
\|W^\top(WCW^\top - B)W\|_F
=
\|C - W^\top B W\|_F
=
\|C - \widetilde{B}\|_F.
\]
Therefore the objective and constraint become
\[
\min_{\widetilde{B}:\ \operatorname{rank}(\widetilde{B})\le k}\ \|C-\widetilde{B}\|_F.
\]
By definition of $C_k$, a minimizer is $\widetilde{B}^\star = C_k$.
Transforming back gives
\[
B^\star = W\widetilde{B}^\star W^\top = WC_kW^\top.
\]
Hence
\[
WC_kW^\top
=
\arg\min_{B:\ \operatorname{rank}(B)\le k}\ \|WCW^\top - B\|_F,
\]
as required.

\item[(b)] Since $A$ has full rank, $\operatorname{rank}(A)=n$. Then $AA^\top\in\mathbb{R}^{n\times n}$ is symmetric positive definite
and in particular $\operatorname{rank}(AA^\top)=n$. The feasible set for ``rank-$n$ approximation'' is
So we can choose $B=AA^\top$, which yields objective value $\|AA^\top-AA^\top\|_F=0$.

\item[(c)] Under the stated assumptions, $R=\operatorname{diag}(r_1,\dots,r_n)$, so
\[
AA^\top = (QR)(QR)^\top = QRR^\top Q^\top
= Q\,\operatorname{diag}(r_1^2,\dots,r_n^2)\,Q^\top.
\]
Thus $AA^\top$ is symmetric PSD with eigen-decomposition $Q\Lambda Q^\top$, where
\[
\Lambda=\operatorname{diag}(r_1^2,\dots,r_n^2),
\qquad
r_1^2>r_2^2>\cdots>r_n^2>0
\]
(using $|r_1|>\cdots>|r_n|$ and $r_i\in\mathbb{R}$).

From part (a), the problem is equivalent to finding the best rank-$k$ approximation to $\Lambda$ in Frobenius norm, and then conjugating by $Q$.

Thus, the best rank-$k$ approximation is obtained by keeping the top $k$ singular values:
\[
(AA^\top)_k = Q S Q^\top,
\qquad
S=\operatorname{diag}(r_1^2,\dots,r_k^2,0,\dots,0).
\]

\item[(d)] In general,
\[
AA^\top = (QR)(QR)^\top = QRR^\top Q^\top.
\]
Let $G$ be the best rank-$k$ approximation to $RR^\top$ in Frobenius norm:
\[
G \in \arg\min_{B:\ \operatorname{rank}(B)\le k}\ \|RR^\top - B\|_F.
\]
Apply part (a) with $C=RR^\top$ and $W=Q$ (orthonormal). Then
\[
QGQ^\top
=
\arg\min_{B:\ \operatorname{rank}(B)\le k}\ \|Q(RR^\top)Q^\top - B\|_F.
\]
Since $Q(RR^\top)Q^\top = AA^\top$, this shows that the best rank-$k$ approximation to $AA^\top$ is $QGQ^\top$.
\end{enumerate}




%----------------------------------------------------------------------------------------
\newpage
\section*{Problem 2. Convexity of Functions}

(a) Non-negative weighted sum.

Show that the non-negative weighted sum of convex functions is convex: i.e. if $f_1, \dots, f_n$ are $n$ convex
functions from $\mathbb{R}^n$ to $\mathbb{R}$ and $w_1, \dots, w_n \in \mathbb{R}_+$ are $n$ positive scalars, then the function:
\[
f = \sum_{i=1}^n w_i f_i \tag{4}
\]
is convex. To make the question easier, you can assume that the functions $f_1, \dots, f_n$ are twice-differentiable.

A set $C$ is convex if and only if the line segment between any two points in $C$ lies in $C$:
\[
C \text{ is convex } \Longleftrightarrow
\forall \vec{x}_1, \vec{x}_2 \in C, \ \forall \theta \in [0,1],
\theta \vec{x}_1 + (1-\theta)\vec{x}_2 \in C. \tag{5}
\]

(b) Prove the convexity of a hyperplane, 
\[
L = \{\vec{x} \mid \vec{a}^\top \vec{x} = b\}.
\]

(c) Prove the convexity of a halfspace,
\[
H = \{\vec{x} \mid \vec{a}^\top \vec{x} \le b\}.
\]

(d) A function $f : \mathbb{R}^n \to \mathbb{R}^m$ is affine if it is the sum of a linear function and a constant,
\[
f(\vec{x}) = A\vec{x} + \vec{b}, \tag{6}
\]
for $A \in \mathbb{R}^{m\times n}$ and $\vec{b} \in \mathbb{R}^m$.

Prove that if $S \subseteq \mathbb{R}^n$ is convex, then the image of $S$ under an affine function $f$,
\[
f(S) = \{ f(\vec{x}) \mid \vec{x} \in S \}, \tag{7}
\]
is convex.

\subsection*{Answer}

\begin{enumerate}
\item[(a)] Since each $f_i$ is convex, for all $x,y \in \mathbb{R}^n$ and $\theta \in [0,1]$,
\[
f_i(\theta x + (1-\theta)y)
\le
\theta f_i(x) + (1-\theta) f_i(y).
\]
\[
w_i f_i(\theta x + (1-\theta)y)
\le
\theta w_i f_i(x) + (1-\theta) w_i f_i(y).
\]
Summing over $i=1,\dots,n$ gives
\[
\sum_{i=1}^n w_i f_i(\theta x + (1-\theta)y)
\le
\theta \sum_{i=1}^n w_i f_i(x)
+
(1-\theta) \sum_{i=1}^n w_i f_i(y).
\]
Thus
\[
f(\theta x + (1-\theta)y)
\le
\theta f(x) + (1-\theta) f(y),
\]
so $f$ is convex.

\item[(b)] Let $\vec{x}_1, \vec{x}_2 \in L$. Then
\[
\vec{a}^\top \vec{x}_1 = b,
\qquad
\vec{a}^\top \vec{x}_2 = b.
\]
For any $\theta \in [0,1]$,
\[
\vec{a}^\top(\theta \vec{x}_1 + (1-\theta)\vec{x}_2)
=
\theta \vec{a}^\top \vec{x}_1
+
(1-\theta)\vec{a}^\top \vec{x}_2
=
\theta b + (1-\theta)b
=
b.
\]
Thus $\theta \vec{x}_1 + (1-\theta)\vec{x}_2 \in L$.
Therefore $L$ is convex.

\item[(c)] Let $\vec{x}_1, \vec{x}_2 \in H$. Then
\[
\vec{a}^\top \vec{x}_1 \le b,
\qquad
\vec{a}^\top \vec{x}_2 \le b.
\]
For $\theta \in [0,1]$,
\[
\vec{a}^\top(\theta \vec{x}_1 + (1-\theta)\vec{x}_2)
=
\theta \vec{a}^\top \vec{x}_1
+
(1-\theta)\vec{a}^\top \vec{x}_2
\le
\theta b + (1-\theta)b
=
b.
\]
Hence $\theta \vec{x}_1 + (1-\theta)\vec{x}_2 \in H$.
Therefore $H$ is convex.

\item[(d)] Let $\vec{x}_1, \vec{x}_2 \in S$.
Since $S$ is convex,
\[
\theta \vec{x}_1 + (1-\theta)\vec{x}_2 \in S
\quad \forall \theta \in [0,1].
\]
Apply $f$:
\begin{align*}
f(\theta \vec{x}_1 + (1-\theta)\vec{x}_2) & = A(\theta \vec{x}_1 + (1-\theta)\vec{x}_2) + \vec{b} \\
 & = \theta (A\vec{x}_1 + \vec{b}) + (1-\theta)(A\vec{x}_2 + \vec{b}) \\
 & = \theta f(\vec{x}_1) + (1-\theta) f(\vec{x}_2).
 \end{align*}
Thus,  for any convex combination of points in $f(S)$, it remains in $f(S)$.
\\$\Rightarrow f(S)$ is convex.
\end{enumerate}

%----------------------------------------------------------------------------------------

\newpage
\section*{Problem 3. Perspective function}

Let $f : \mathbb{R}^n \to \mathbb{R}$ be a differentiable convex function, with domain the entire space $\mathbb{R}^n$.

(a) Show that we can represent $f$ as a pointwise maximum of affine functions, specifically
\[
\forall x \in \mathbb{R}^n :
f(x) = \max_{(a,b)\in \mathcal{A}} a^\top x + b, \tag{8}
\]
where $\mathcal{A} \subseteq \mathbb{R}^{n+1}$ is to be determined.

Hint: For every $x, y \in \mathbb{R}^n$:
\[
f(x) \ge f(y) + (x-y)^\top \nabla f(y).
\]

(b) We define the perspective of $f$ as the function $g : \mathbb{R}^n \times \mathbb{R} \to \mathbb{R}$ with values
\[
g(x,t) =
\begin{cases}
t f(x/t) & \text{if } t>0, \\
+\infty & \text{otherwise}.
\end{cases}
\]

Prove that $g$ is convex.

Hint: Use (8).

(c) Show that the function $h : \mathbb{R}^n \to \mathbb{R}$ with values
\[
h(x) =
\begin{cases}
\frac{(a^\top x + b)^2}{c^\top x + d} & \text{if } c^\top x + d > 0, \\
+\infty & \text{otherwise},
\end{cases}
\]
is convex.

Hint: Consider the perspective function of $f(z)=z^2$ on $\mathbb{R}$.

\subsection*{Answer}
\begin{enumerate}
\item[(a)] For any fixed $y\in\mathbb{R}^n$, define the affine function $\ell_y:\mathbb{R}^n\to\mathbb{R}$ by
\[
\ell_y(x) \doteq f(y) + \nabla f(y)^\top (x-y)
= (\nabla f(y))^\top x + \big(f(y)-\nabla f(y)^\top y\big).
\]
By the first-order condition for differentiable convex functions (supporting hyperplane inequality),
\[
f(x)\ge \ell_y(x)\qquad \forall x\in\mathbb{R}^n,\ \forall y\in\mathbb{R}^n. \tag{*}
\]
In particular, for each fixed $x$,
\[
f(x)\ge \sup_{y\in\mathbb{R}^n} \ell_y(x).
\]
On the other hand, choosing $y=x$ gives $\ell_x(x)=f(x)$, so
\[
\sup_{y\in\mathbb{R}^n}\ell_y(x)\ge \ell_x(x)=f(x).
\]
Thus
\[
f(x)=\sup_{y\in\mathbb{R}^n}\ell_y(x).
\]
Now define
\[
A \doteq \left\{(a,b)\in\mathbb{R}^{n+1}:\ \exists\,y\in\mathbb{R}^n \text{ such that } a=\nabla f(y),\ b=f(y)-\nabla f(y)^\top y\right\}.
\]
Then $\ell_y(x)=a^\top x+b$ for $(a,b)$ generated by $y$, and the previous identity becomes
\[
f(x)=\max_{(a,b)\in A} a^\top x+b,
\]
where $\max$ may be interpreted as $\sup$ if the maximum is not attained. This proves (8).

\item[(b)] Using (8), for $t>0$ we have
\begin{align*}
g(x,t)
&= t f(x/t)
= t \max_{(a,b)\in A}\left(a^\top (x/t)+b\right) \\
&= \max_{(a,b)\in A}\left( a^\top x + bt \right).
\end{align*}
For each fixed $(a,b)\in A$, the map $(x,t)\mapsto a^\top x+bt$ is affine on $\mathbb{R}^n\times\mathbb{R}$.
A pointwise maximum of affine functions is convex, hence the function
\[
(x,t)\mapsto \max_{(a,b)\in A}(a^\top x+bt)
\]
is convex on $\{(x,t):t>0\}$. Extending $g$ by $+\infty$ on $\{t\le 0\}$ preserves convexity (equivalently, $\operatorname{epi}(g)$ remains convex),
so $g$ is convex on $\mathbb{R}^n\times\mathbb{R}$.

\item[(c)] Consider $f:\mathbb{R}\to\mathbb{R}$ given by $f(z)=z^2$, which is convex and differentiable on $\mathbb{R}$.
By part (b), its perspective $p:\mathbb{R}\times\mathbb{R}\to\mathbb{R}$ defined by
\[
p(z,t)=
\begin{cases}
t f(z/t) & t>0,\\
+\infty & \text{otherwise},
\end{cases}
\]
is convex. For $t>0$,
\[
p(z,t)=t\left(\frac{z}{t}\right)^2=\frac{z^2}{t}.
\]
Now define the affine map $\phi:\mathbb{R}^n\to\mathbb{R}\times\mathbb{R}$ by
\[
\phi(x) \doteq (z(x),t(x)) = (a^\top x+b,\ c^\top x+d).
\]
Then $h$ can be written as the composition
\[
h(x)=p(\phi(x))=p(a^\top x+b,\ c^\top x+d).
\]
Since $p$ is convex and $\phi$ is affine, the composition $p\circ\phi$ is convex. Writing it out yields exactly
\[
h(x)=
\begin{cases}
\dfrac{(a^\top x+b)^2}{c^\top x+d} & \text{if } c^\top x+d>0,\\[6pt]
+\infty & \text{otherwise},
\end{cases}
\]
so $h$ is convex.
\end{enumerate}

%----------------------------------------------------------------------------------------
\newpage
\section*{Problem 4. Convexity of Sets}

Determine if each set $C$ given below is convex. Prove that each set is convex or provide an example to show that
it is not convex. You may use any techniques used in class or discussion to demonstrate or disprove convexity.

(a) $C = \{\vec{x} \in \mathbb{R}^2 \mid x_1 x_2 \ge 0\}$, where $\vec{x} = [x_1 \ x_2]^\top$.

(b) $C = \{X \in S^n \mid \lambda_{\min}(X) \ge 2\}$, where $S^n$ is the set of symmetric matrices in $\mathbb{R}^{n\times n}$ and $\lambda_{\min}(X)$ is the minimum eigenvalue of $X$.

(c) Let $H(\vec{w})$ denote the hyperplane with normal direction $\vec{w} \in \mathbb{R}^n$, i.e.,
\[
H(\vec{w}) = \{\vec{x} \in \mathbb{R}^n \mid \vec{x}^\top \vec{w} = 0\}. \tag{9}
\]

Let $P : \mathbb{R}^n \to \mathbb{R}^n$ be given by
\[
P(\vec{x}) = \arg\min_{\vec{y} \in H(\vec{w})} \|\vec{y} - \vec{x}\|_2. \tag{10}
\]

Let
\[
C = \{ P(\vec{x}) \mid \vec{x} \in B \}, \tag{11}
\]
where $B = \{\vec{x} \in \mathbb{R}^n \mid \|\vec{x}\|_2 \le 1\}$.

\subsection*{Answer}

\begin{enumerate}
\item[(a)] The set $C$ is not convex.
\[
\vec{x}_1=\begin{bmatrix}1\\[2pt]0\end{bmatrix}\in C
\qquad
\vec{x}_2=\begin{bmatrix}0\\[2pt]-1\end{bmatrix}\in C
\]
For $\theta=\tfrac12$,
\[
\vec{x}_\theta=\frac12\vec{x}_1+\frac12\vec{x}_2
=\begin{bmatrix}\tfrac12\\[2pt]-\tfrac12\end{bmatrix},
\qquad
(\vec{x}_\theta)_1(\vec{x}_\theta)_2
=
\left(\tfrac12\right)\left(-\tfrac12\right)
=
-\tfrac14
<0,
\]
so $\vec{x}_\theta\notin C$. Hence $C$ is not convex.

\item[(b)] The set $C$ is convex.

Let $X_1,X_2\in C$ and $\theta\in[0,1]$, and define $X_\theta=\theta X_1+(1-\theta)X_2$.
Using the Rayleigh quotient characterization,
\[
\lambda_{\min}(X)=\min_{\|v\|_2=1} v^\top X v.
\]
Fix any $v$ with $\|v\|_2=1$. Since $X_1,X_2\in C$, we have
\[
v^\top X_1 v \ge \lambda_{\min}(X_1)\ge 2,
\qquad
v^\top X_2 v \ge \lambda_{\min}(X_2)\ge 2.
\]
Therefore
\[
v^\top X_\theta v
=
\theta v^\top X_1 v + (1-\theta) v^\top X_2 v
\ge
\theta\cdot 2 + (1-\theta)\cdot 2
=
2.
\]
Since this holds for all unit $v$,
\[
\lambda_{\min}(X_\theta)
=
\min_{\|v\|_2=1} v^\top X_\theta v
\ge 2,
\]
so $X_\theta\in C$. Hence $C$ is convex.

\item[(c)] The set $C$ is convex.

First, note that $P(\vec{x})$ is the Euclidean projection of $\vec{x}$ onto the hyperplane $H(\vec{w})$.
We can compute it explicitly. Let $\vec{w}\neq \vec{0}$ and write $\vec{y}=\vec{x}-\alpha \vec{w}$.
The constraint $\vec{y}^\top \vec{w}=0$ gives
\[
(\vec{x}-\alpha \vec{w})^\top \vec{w}=0
\quad\Longrightarrow\quad
\vec{x}^\top \vec{w}-\alpha \|\vec{w}\|_2^2=0
\quad\Longrightarrow\quad
\alpha=\frac{\vec{x}^\top \vec{w}}{\|\vec{w}\|_2^2}.
\]
Thus
\[
P(\vec{x})
=
\vec{x}-\frac{\vec{x}^\top \vec{w}}{\|\vec{w}\|_2^2}\,\vec{w}
=
\left(I-\frac{\vec{w}\vec{w}^\top}{\|\vec{w}\|_2^2}\right)\vec{x}.
\]
Hence $P$ is linear (and therefore affine).

Since $B=\{\vec{x}\in\mathbb{R}^n:\|\vec{x}\|_2\le 1\}$ is convex, and the image of a convex set under an affine map is convex.
\end{enumerate}

%----------------------------------------------------------------------------------------
\newpage
\section*{Problem 5. PSD Matrices}

In this problem, we will analyze properties of positive semidefinite (PSD) matrices. A symmetric matrix
$M \in \mathbb{R}^{n\times n}$ is a PSD matrix if $\vec{x}^\top M \vec{x} \ge 0$ for all $\vec{x} \in \mathbb{R}^n$, and we denote that as
$M \succeq 0$ or $M \in S^n_+$.

Assume $A \in \mathbb{R}^{n\times n}$ is a symmetric matrix.

(a) Show that $A \succeq 0$ if and only if all eigenvalues of $A$ are non-negative.

(b) Show that if $A \succeq 0$ then all diagonal entries of $A$ are non-negative, $A_{ii} \ge 0$.

Now we will show that $A \in S^n_+$ (i.e., $A \in \mathbb{R}^{n\times n}$ and $A \succeq 0$) if and only if there exists $P \in S^n_+$ such that
\[
A = P^\top P = P^2.
\]
Such a matrix $P$ is known as a PSD square root of $A$.

(c) First, show that if $A \in S^n_+$, then there exists $P \in S^n_+$ such that $A = P^2$.

(d) Now, show that for any matrix $Q \in \mathbb{R}^{m\times n}$, if $A = Q^\top Q$ then $A \in S^n_+$.

NOTE: If we take $Q = P \in S^n_+$, we have $A = P^2 \in S^n_+$; this proves the other direction of the above
“if-and-only-if.”

We will use positive semidefinite matrices to construct the singular value decomposition (SVD). One can construct
the SVD of a matrix $B \in \mathbb{R}^{m\times n}$ in two ways, using either the matrices $BB^\top \in \mathbb{R}^{m\times m}$ or $B^\top B \in \mathbb{R}^{n\times n}$, and
usually one chooses the method depending on which matrix is smaller. The following result shows that the
singular values will be the same no matter what construction you use.

(e) Let $B \in \mathbb{R}^{m\times n}$ be an arbitrary matrix (not necessarily PSD or even square). From the previous parts of
this problem, $BB^\top$ and $B^\top B$ are PSD, thus having real non-negative eigenvalues. Prove that the non-zero
eigenvalues of $BB^\top$ are the same as the non-zero eigenvalues of $B^\top B$.


\subsection*{Answer}

\begin{enumerate}
\item[(a)] Since $A$ is symmetric, it has an orthonormal eigen-decomposition
\[
A = U\Lambda U^\top,
\]
where $U$ is orthonormal and $\Lambda=\operatorname{diag}(\lambda_1,\dots,\lambda_n)$ contains the real eigenvalues of $A$.

\emph{($\Rightarrow$)} Suppose $A\succeq 0$. Let $\lambda_i$ be an eigenvalue with eigenvector $\vec{u}_i$ (the $i$th column of $U$),
so $A\vec{u}_i=\lambda_i \vec{u}_i$ and $\|\vec{u}_i\|_2=1$. Then
\[
0 \le \vec{u}_i^\top A \vec{u}_i = \vec{u}_i^\top (\lambda_i \vec{u}_i)=\lambda_i \|\vec{u}_i\|_2^2=\lambda_i,
\]
so $\lambda_i\ge 0$ for all $i$.

\emph{($\Leftarrow$)} Conversely, suppose all eigenvalues satisfy $\lambda_i\ge 0$. For any $\vec{x}\in\mathbb{R}^n$, write $\vec{x}=U\vec{z}$ where
$\vec{z}=U^\top \vec{x}$. Then
\[
\vec{x}^\top A \vec{x}
=
(U\vec{z})^\top (U\Lambda U^\top)(U\vec{z})
=
\vec{z}^\top \Lambda \vec{z}
=
\sum_{i=1}^n \lambda_i z_i^2
\ge 0,
\]
since $\lambda_i\ge 0$ and $z_i^2\ge 0$. Hence $A\succeq 0$.

\item[(b)] Let $\vec{e}_i$ be the $i$th standard basis vector. If $A\succeq 0$, then for all $i$,
\[
0 \le \vec{e}_i^\top A \vec{e}_i = A_{ii}.
\]
Thus $A_{ii}\ge 0$ for all $i$.

\item[(c)] Assume $A\in S^n_+$, i.e. $A$ is symmetric and PSD. By part (a), all eigenvalues of $A$ are non-negative.
Let
\[
A = U\Lambda U^\top,
\qquad
\Lambda=\operatorname{diag}(\lambda_1,\dots,\lambda_n),\ \lambda_i\ge 0.
\]
Define
\[
\Lambda^{1/2} \doteq \operatorname{diag}(\sqrt{\lambda_1},\dots,\sqrt{\lambda_n}),
\qquad
P \doteq U\Lambda^{1/2}U^\top.
\]
Then $P$ is symmetric. Moreover,
\[
P^2
=
(U\Lambda^{1/2}U^\top)(U\Lambda^{1/2}U^\top)
=
U\Lambda^{1/2}(U^\top U)\Lambda^{1/2}U^\top
=
U\Lambda U^\top
=
A.
\]
Finally, $P\succeq 0$ because its eigenvalues are $\sqrt{\lambda_i}\ge 0$ (again by part (a) applied to $P$).
Thus there exists $P\in S^n_+$ such that $A=P^2$.

\item[(d)] Let $Q\in\mathbb{R}^{m\times n}$ and define $A=Q^\top Q$. Then $A$ is symmetric since
\[
A^\top = (Q^\top Q)^\top = Q^\top Q = A.
\]
For any $\vec{x}\in\mathbb{R}^n$,
\[
\vec{x}^\top A \vec{x}
=
\vec{x}^\top Q^\top Q \vec{x}
=
(Q\vec{x})^\top (Q\vec{x})
=
\|Q\vec{x}\|_2^2
\ge 0.
\]
Hence $A\succeq 0$, i.e. $A\in S^n_+$.

\item[(e)] Let $B\in\mathbb{R}^{m\times n}$ be arbitrary. From (d), both $BB^\top$ and $B^\top B$ are PSD and hence have real non-negative eigenvalues.

We show that the \emph{non-zero} eigenvalues of $BB^\top$ and $B^\top B$ coincide (including algebraic multiplicity).

First, let $\lambda>0$ and suppose $\lambda$ is an eigenvalue of $B^\top B$ with eigenvector $\vec{v}\neq \vec{0}$:
\[
B^\top B \vec{v} = \lambda \vec{v}. \tag{$\ast$}
\]
Then $B\vec{v}\neq \vec{0}$; otherwise if $B\vec{v}=\vec{0}$, multiplying by $B^\top$ gives $B^\top B\vec{v}=\vec{0}$, contradicting $\lambda>0$ and $\vec{v}\neq 0$.
Multiply ($\ast$) on the left by $B$:
\[
BB^\top (B\vec{v}) = B(B^\top B\vec{v}) = B(\lambda \vec{v}) = \lambda (B\vec{v}).
\]
Thus $\lambda$ is an eigenvalue of $BB^\top$ with eigenvector $B\vec{v}\neq 0$.

Conversely, let $\lambda>0$ and suppose $\lambda$ is an eigenvalue of $BB^\top$ with eigenvector $\vec{u}\neq \vec{0}$:
\[
BB^\top \vec{u} = \lambda \vec{u}. \tag{$\dagger$}
\]
Then $B^\top \vec{u}\neq \vec{0}$ (same reasoning: if $B^\top \vec{u}=0$, then $BB^\top \vec{u}=0$, contradicting $\lambda>0$).
Multiply ($\dagger$) on the left by $B^\top$:
\[
B^\top B (B^\top \vec{u}) = B^\top (BB^\top \vec{u}) = B^\top (\lambda \vec{u}) = \lambda (B^\top \vec{u}).
\]
Thus $\lambda$ is an eigenvalue of $B^\top B$ with eigenvector $B^\top \vec{u}\neq 0$.

Therefore, $\lambda>0$ is an eigenvalue of $B^\top B$ if and only if it is an eigenvalue of $BB^\top$.

To match multiplicities: 
for real symmetric matrices, the algebraic and geometric multiplicities coincide, so the eigenspaces $E_\lambda(B^\top B)$ and $E_\lambda(BB^\top)$ have dimensions equal to the multiplicity of $\lambda$ as an eigenvalue of $B^\top B$ and $BB^\top$, respectively. 

Hence $\dim E_\lambda(B^\top B)=\dim E_\lambda(BB^\top)$ for every $\lambda>0$, so the non-zero eigenvalues coincide with the same multiplicities.
\end{enumerate}
%----------------------------------------------------------------------------------------

\newpage
\section*{Problem 6. Homework Process}

With whom did you work on this homework? List the names and SIDs of your group members.

NOTE: If you didn’t work with anyone, you can put “none” as your answer.

\subsection*{Answer}
none

%----------------------------------------------------------------------------------------
\end{document}
